\section{Exercise 1a: Coherent Demodulation}
\label{sec:lab2_ex1a}

\subsection{Objective}

Implement a coherent AM demodulator that recovers the baseband message in both
time and frequency domains.

\subsection{Module Design}

The coherent demodulation sub-VI implements the following processing chain:
\begin{enumerate}
\item multiply the AM input signal by a local carrier at the same carrier
frequency,
\item apply a Butterworth low-pass filter with adjustable order and cutoff,
\item estimate the DC component of the filtered signal and subtract it,
\item display the recovered message in time domain and by means of a PSD
estimate.
\end{enumerate}

\begin{figure}[H]
\centering
\labincludegraphics[width=\linewidth]{figures/Coherant_Demodulation_VI.png}
\caption{Block diagram of the coherent demodulation VI.}
\label{fig:lab2_coherent_vi}
\end{figure}

\subsection{Link to Theory}

From the coherent detection model in
Section~\ref{sec:lab2_theory}, the low-pass filtered output is
\begin{equation}
y_{\mathrm{LPF}}(t)=\frac{1}{2}\left[A_c + A_m m(t)\right]\cos\phi.
\end{equation}
This shows that the demodulator output contains both the desired baseband term
and a DC offset proportional to the carrier amplitude.

\subsection{Why Filtering, DC Removal, and Scaling Are Needed}

\begin{itemize}
\item \textbf{Low-pass filter:} removes the mixing product around $2f_c$.
\item \textbf{DC removal:} removes the carrier-dependent offset and recenters the
output around zero.
\item \textbf{Amplitude scaling:} compensates for the attenuation introduced by
the product detector.
\end{itemize}

\subsection{Expected Output}

With correct carrier synchronization, the recovered waveform should be a sinusoid
at the message frequency $f_m$. In the PSD, the dominant spectral component
should also appear at $f_m$, confirming successful baseband recovery.
